% Python Language Reference Card
\documentclass[9pt]{extarticle}
\usepackage[landscape,paperwidth=11in,paperheight=8.5in,
            margin=0.4in]{geometry}
\usepackage{tabularx}
\usepackage{multicol}
\usepackage{needspace}
\usepackage[T1]{fontenc}
\usepackage{microtype}
\usepackage{helvet}
\usepackage[scaled=0.82]{beramono}
\renewcommand{\familydefault}{\sfdefault}
\newcolumntype{Y}{>{\raggedright\arraybackslash}X}

\pagestyle{empty}
\setlength{\parindent}{0pt}
\setlength{\parskip}{0pt}
\setlength{\columnsep}{0.3in}

\newcommand{\cardtitle}[1]{%
  {\Large\bfseries #1}\par\vspace{1pt}%
  \rule{\linewidth}{0.4pt}\par\vspace{4pt}%
}
\newcommand{\sect}[1]{\needspace{4\baselineskip}\vspace{6pt}{\footnotesize\bfseries\scshape #1}\par\vspace{2pt}}

\begin{document}\small

\cardtitle{Python}

% ── Standard Tasks ─────────────────────────────────────────────────
\begin{tabularx}{\linewidth}{@{}l@{\quad}l@{\quad}Y@{}}
\textbf{Task}        & \textbf{Command}                                  & \textbf{Notes} \\[2pt]
Version manager      & \texttt{uv}                                        & \texttt{uv python install 3.12} to install a version \\
New project          & \texttt{uv init \textit{name}}                     & \texttt{-{}-lib} for library; creates \texttt{pyproject.toml} \\
Add dependency       & \texttt{uv add \textit{pkg}}                       & \texttt{-{}-dev} for dev-only dependency \\
Install deps         & \texttt{uv sync}                                   & installs from \texttt{uv.lock} \\
Build                & \texttt{uv build}                                  & builds wheel + sdist to \texttt{dist/} \\
Run                  & \texttt{uv run \textit{script}.py}                 & runs in project virtualenv \\
REPL                 & \texttt{uv run python} / \texttt{ipython}          & \\
Test                 & \texttt{uv run pytest}                             & \texttt{-x} stop first fail, \texttt{-k \textit{pat}} filter \\
Format               & \texttt{ruff format .}                             & \texttt{-{}-check} for dry run \\
Lint                 & \texttt{ruff check .}                              & \texttt{-{}-fix} to autofix \\
LSP                  & \texttt{pyright}                                   & configured in neovim via lspconfig \\
\end{tabularx}

\vspace{6pt}
\rule{\linewidth}{0.2pt}

% ── Language-Specific Content ──────────────────────────────────────
\begin{multicols}{3}

\sect{uv Extras}
\begin{tabularx}{\linewidth}{@{}l@{\enspace}Y@{}}
\texttt{uv remove \textit{pkg}}       & Remove dependency \\
\texttt{uv lock}                       & Update lockfile \\
\texttt{uv tool install \textit{pkg}} & Install global CLI \\
\texttt{uv tool list}                  & List global tools \\
\texttt{uvx \textit{cmd}}             & Run tool without install \\
\texttt{uv python list}               & Available versions \\
\texttt{uv python pin 3.12}           & Pin project version \\
\end{tabularx}

\sect{Common Patterns}
\begin{tabularx}{\linewidth}{@{}l@{\enspace}Y@{}}
\texttt{[x for x in xs if p(x)]}     & List comprehension \\
\texttt{\{k: v for k, v in d\}}      & Dict comprehension \\
\texttt{x if cond else y}            & Ternary \\
\texttt{a, b = b, a}                 & Swap \\
\texttt{*args, **kwargs}             & Variadic args \\
\texttt{yield}                       & Generator \\
\texttt{with open(f) as fh:}         & Context manager \\
\texttt{match x:}                    & Pattern match (3.10+) \\
\texttt{case \textit{pat}:}          & Match arm \\
\texttt{f"\{x:.2f\}"}               & f-string formatting \\
\end{tabularx}

\sect{Error Handling}
\begin{tabularx}{\linewidth}{@{}l@{\enspace}Y@{}}
\texttt{try: ... except E as e:}     & Catch exception \\
\texttt{except (A, B):}              & Catch multiple types \\
\texttt{else:}                       & No exception branch \\
\texttt{finally:}                    & Always runs \\
\texttt{raise ValueError("msg")}     & Raise exception \\
\texttt{raise X from e}              & Chain exceptions \\
\texttt{class E(Exception): ...}     & Custom exception \\
\end{tabularx}

\sect{Strings \& Slicing}
\begin{tabularx}{\linewidth}{@{}l@{\enspace}Y@{}}
\texttt{" ".join(xs)}                & Join with separator \\
\texttt{s.split(",")  }             & Split string \\
\texttt{s.strip()}                   & Trim whitespace \\
\texttt{s.startswith("x")}          & Prefix test \\
\texttt{s.replace(a, b)}            & Replace substring \\
\texttt{xs[1:-1]}                    & Slice (drop first/last) \\
\texttt{xs[::2]}                     & Every 2nd element \\
\texttt{xs[::-1]}                    & Reverse \\
\end{tabularx}

\sect{Type Hints}
\begin{tabularx}{\linewidth}{@{}l@{\enspace}Y@{}}
\texttt{x: int = 0}               & Variable annotation \\
\texttt{def f(x: str) -> bool:}   & Function signature \\
\texttt{list[int]}                 & Generic list \\
\texttt{dict[str, Any]}           & Generic dict \\
\texttt{X | None}                  & Optional (3.10+) \\
\texttt{TypeAlias = ...}           & Type alias \\
\texttt{@dataclass}                & Data class \\
\texttt{Protocol}                  & Structural typing \\
\texttt{\# pyright: strict}        & Strict mode (file) \\
\end{tabularx}

\sect{Standard Library}
\begin{tabularx}{\linewidth}{@{}l@{\enspace}Y@{}}
\texttt{pathlib}          & File paths \\
\texttt{itertools}        & chain, product, groupby \\
\texttt{functools}        & lru\_cache, partial, reduce \\
\texttt{collections}      & Counter, defaultdict, deque \\
\texttt{dataclasses}      & @dataclass \\
\texttt{typing}           & Type hint utilities \\
\texttt{json}             & JSON encode/decode \\
\texttt{subprocess}       & Run shell commands \\
\texttt{argparse}         & CLI argument parsing \\
\texttt{logging}          & Structured logging \\
\end{tabularx}

\sect{Key Libraries}
\begin{tabularx}{\linewidth}{@{}l@{\enspace}Y@{}}
\texttt{polars}           & Fast DataFrames \\
\texttt{numpy}            & Array computing \\
\texttt{scipy}            & Numerics, optimization \\
\texttt{matplotlib}       & Plotting \\
\texttt{xarray}           & N-dim labeled arrays \\
\texttt{geopandas}        & Geospatial DataFrames \\
\texttt{httpx}            & HTTP client (async) \\
\texttt{pydantic}         & Validation / settings \\
\texttt{pytest}           & Testing framework \\
\end{tabularx}

\sect{Decorators}
\begin{tabularx}{\linewidth}{@{}l@{\enspace}Y@{}}
\texttt{@property}                    & Getter method \\
\texttt{@x.setter}                    & Setter method \\
\texttt{@staticmethod}                & No self/cls arg \\
\texttt{@classmethod}                 & Receives cls arg \\
\texttt{@abstractmethod}              & Must override (ABC) \\
\texttt{def dec(f): ... return w}     & Custom decorator \\
\texttt{@wraps(f)}                    & Preserve metadata \\
\end{tabularx}

\sect{Async}
\begin{tabularx}{\linewidth}{@{}l@{\enspace}Y@{}}
\texttt{async def f():}               & Async function \\
\texttt{await coro}                   & Await coroutine \\
\texttt{asyncio.run(main())}          & Run event loop \\
\texttt{asyncio.gather(*coros)}       & Run concurrently \\
\texttt{async with}                   & Async context manager \\
\texttt{async for x in ait:}          & Async iteration \\
\end{tabularx}

\sect{Data Tools}
\begin{tabularx}{\linewidth}{@{}l@{\enspace}Y@{}}
\texttt{duckdb}                        & SQL REPL (parquet/CSV) \\
\texttt{vd \textit{file}}             & VisiData TUI \\
\texttt{marimo edit \textit{nb}.py}   & Reactive notebook \\
\texttt{quarto render \textit{f}.qmd} & Render document \\
\end{tabularx}

\sect{Testing Patterns}
\begin{tabularx}{\linewidth}{@{}l@{\enspace}Y@{}}
\texttt{def test\_\textit{name}():}    & Test function \\
\texttt{assert x == y}                & Assertion \\
\texttt{@pytest.fixture}              & Fixture decorator \\
\texttt{@pytest.mark.parametrize}     & Parametrize tests \\
\texttt{pytest.raises(Exc)}           & Expect exception \\
\texttt{tmp\_path}                    & Built-in temp dir fixture \\
\end{tabularx}

\sect{Functional Patterns}
\begin{tabularx}{\linewidth}{@{}l@{\enspace}Y@{}}
\texttt{map(f, xs)}                   & Apply f to each element \\
\texttt{filter(f, xs)}                & Keep elements where f is true \\
\texttt{functools.reduce(f, xs)}      & Fold / accumulate \\
\texttt{functools.partial(f, a)}      & Partial application \\
\texttt{@functools.lru\_cache}        & Memoization \\
\texttt{lambda x: \textit{expr}}      & Anonymous function \\
\texttt{zip(xs, ys)}                  & Pair elements \\
\texttt{enumerate(xs)}                & Index + element \\
\texttt{any(p(x) for x in xs)}       & Existential \\
\texttt{all(p(x) for x in xs)}       & Universal \\
\texttt{sorted(xs, key=f)}           & Sort by projection \\
\texttt{itertools.chain(a, b)}        & Concatenate iterables \\
\texttt{itertools.product(a, b)}      & Cartesian product \\
\texttt{itertools.starmap(f, pairs)}  & Map unpacking args \\
\texttt{more-itertools}               & Extended combinators (pip) \\
\end{tabularx}

\sect{FP Libraries}
\begin{tabularx}{\linewidth}{@{}l@{\enspace}Y@{}}
\texttt{toolz / cytoolz}     & Composable FP: pipe, curry, merge \\
\texttt{returns}              & Result, Maybe, IO monads \\
\texttt{expression}           & Pipe, effect system, Result \\
\texttt{more-itertools}       & 200+ iterator recipes \\
\end{tabularx}

\end{multicols}

\newpage

\cardtitle{Python --- Scientific \& Geospatial}

\begin{multicols}{3}

\sect{NumPy --- Creation}
\begin{tabularx}{\linewidth}{@{}l@{\enspace}Y@{}}
\texttt{np.array([1, 2, 3])}        & From list \\
\texttt{np.zeros((m, n))}           & Zero matrix \\
\texttt{np.ones((m, n))}            & Ones matrix \\
\texttt{np.eye(n)}                   & Identity matrix \\
\texttt{np.arange(0, 10, 0.5)}      & Range with step \\
\texttt{np.linspace(0, 1, 100)}     & Evenly spaced \\
\texttt{np.random.default\_rng()}   & RNG instance \\
\texttt{rng.normal(0, 1, (m, n))}   & Normal samples \\
\texttt{np.full((m, n), v)}         & Fill with value \\
\end{tabularx}

\sect{NumPy --- Indexing \& Shape}
\begin{tabularx}{\linewidth}{@{}l@{\enspace}Y@{}}
\texttt{a[1, 2]}                     & Element access \\
\texttt{a[1:3, :]}                   & Row slice \\
\texttt{a[:, 0]}                     & Column slice \\
\texttt{a[a > 0]}                    & Boolean indexing \\
\texttt{a.shape}                     & Dimensions tuple \\
\texttt{a.reshape(m, n)}             & Reshape \\
\texttt{a.T}                         & Transpose \\
\texttt{a.flatten()}                 & 1-D copy \\
\texttt{np.concatenate([a, b])}     & Join arrays \\
\texttt{np.stack([a, b])}           & Stack new axis \\
\end{tabularx}

\sect{NumPy --- Operations}
\begin{tabularx}{\linewidth}{@{}l@{\enspace}Y@{}}
\texttt{a + b, a * b}               & Elementwise ops \\
\texttt{a @ b}                       & Matrix multiply \\
\texttt{np.dot(a, b)}               & Dot product \\
\texttt{a.sum(axis=0)}              & Sum along axis \\
\texttt{a.mean(), a.std()}           & Statistics \\
\texttt{a.min(), a.max()}            & Min / max \\
\texttt{np.argmax(a)}                & Index of max \\
\texttt{np.where(cond, x, y)}       & Conditional select \\
\texttt{np.sort(a)}                  & Sort array \\
\texttt{np.unique(a)}                & Unique values \\
\end{tabularx}

\sect{NumPy --- Linear Algebra}
\begin{tabularx}{\linewidth}{@{}l@{\enspace}Y@{}}
\texttt{np.linalg.inv(A)}           & Inverse \\
\texttt{np.linalg.det(A)}           & Determinant \\
\texttt{np.linalg.eig(A)}           & Eigenvalues/vectors \\
\texttt{np.linalg.svd(A)}           & SVD \\
\texttt{np.linalg.solve(A, b)}      & Solve Ax = b \\
\texttt{np.linalg.norm(x)}          & Vector/matrix norm \\
\texttt{np.linalg.lstsq(A, b)}      & Least squares \\
\end{tabularx}

\sect{Pandas --- DataFrame}
\begin{tabularx}{\linewidth}{@{}l@{\enspace}Y@{}}
\texttt{pd.read\_csv(f)}            & Read CSV \\
\texttt{pd.read\_parquet(f)}        & Read Parquet \\
\texttt{df.to\_csv(f)}              & Write CSV \\
\texttt{df.head(n)}                  & First n rows \\
\texttt{df.info()}                   & Column types + nulls \\
\texttt{df.describe()}               & Summary stats \\
\texttt{df.shape}                    & (rows, cols) \\
\texttt{df.columns}                  & Column names \\
\texttt{df.dtypes}                   & Column types \\
\end{tabularx}

\sect{Pandas --- Selection}
\begin{tabularx}{\linewidth}{@{}l@{\enspace}Y@{}}
\texttt{df["col"]}                   & Column (Series) \\
\texttt{df[["a", "b"]]}             & Multiple columns \\
\texttt{df.loc[row, col]}           & Label-based \\
\texttt{df.iloc[i, j]}              & Position-based \\
\texttt{df[df.x > 0]}               & Boolean filter \\
\texttt{df.query("x > 0")}          & Query string \\
\texttt{df.sort\_values("col")}     & Sort by column \\
\texttt{df.drop\_duplicates()}      & Remove duplicates \\
\texttt{df.isna()}                   & Null check \\
\texttt{df.fillna(v)}               & Fill nulls \\
\end{tabularx}

\sect{Pandas --- Transform}
\begin{tabularx}{\linewidth}{@{}l@{\enspace}Y@{}}
\texttt{df.groupby("k").agg(f)}     & Group + aggregate \\
\texttt{df.pivot\_table(...)}        & Pivot / crosstab \\
\texttt{pd.merge(a, b, on="k")}     & Join DataFrames \\
\texttt{pd.concat([a, b])}          & Stack DataFrames \\
\texttt{df.apply(f)}                 & Apply function \\
\texttt{df.assign(y=lambda d: ...)} & Add column \\
\texttt{df.rename(columns=\{...\})} & Rename columns \\
\texttt{df.melt(...)}                & Wide to long \\
\end{tabularx}

\sect{Math Module}
\begin{tabularx}{\linewidth}{@{}l@{\enspace}Y@{}}
\texttt{math.sqrt(x)}               & Square root \\
\texttt{math.sin, cos, tan}         & Trig functions \\
\texttt{math.log(x)} / \texttt{log10}  & Ln / log base 10 \\
\texttt{math.exp(x)}                & $e^x$ \\
\texttt{math.pi, math.e}            & Constants \\
\texttt{math.ceil, math.floor}      & Round up / down \\
\texttt{math.gcd(a, b)}             & GCD \\
\texttt{math.comb(n, k)}            & Binomial coeff \\
\texttt{math.isclose(a, b)}         & Float comparison \\
\end{tabularx}

\sect{Shapely --- Geometry}
\begin{tabularx}{\linewidth}{@{}l@{\enspace}Y@{}}
\texttt{Point(x, y)}                & Point geometry \\
\texttt{LineString([(x,y), ...])}   & Line geometry \\
\texttt{Polygon([(x,y), ...])}      & Polygon geometry \\
\texttt{g.area, g.length}           & Measurements \\
\texttt{g.buffer(d)}                & Buffer distance \\
\texttt{g.intersection(h)}          & Intersection \\
\texttt{g.union(h)}                 & Union \\
\texttt{g.contains(h)}              & Containment test \\
\texttt{g.distance(h)}              & Distance between \\
\texttt{g.centroid}                  & Centroid point \\
\end{tabularx}

\sect{GeoPandas}
\begin{tabularx}{\linewidth}{@{}l@{\enspace}Y@{}}
\texttt{gpd.read\_file(f)}          & Read shapefile/gpkg \\
\texttt{gpd.read\_parquet(f)}       & Read GeoParquet \\
\texttt{gdf.to\_file(f)}            & Write to file \\
\texttt{gdf.geometry}               & Geometry column \\
\texttt{gdf.crs}                     & Coordinate ref system \\
\texttt{gdf.to\_crs(epsg=4326)}     & Reproject \\
\texttt{gdf.plot()}                  & Quick plot \\
\texttt{gpd.sjoin(a, b)}            & Spatial join \\
\texttt{gdf.dissolve(by="col")}     & Merge geometries \\
\texttt{gdf.overlay(b, how="...")} & Spatial overlay \\
\texttt{gdf.clip(mask)}             & Clip to boundary \\
\end{tabularx}

\sect{Rasterio}
\begin{tabularx}{\linewidth}{@{}l@{\enspace}Y@{}}
\texttt{rasterio.open(f)}           & Open raster file \\
\texttt{src.read(1)}                & Read band 1 (array) \\
\texttt{src.bounds}                  & Spatial extent \\
\texttt{src.crs}                     & Coordinate ref system \\
\texttt{src.transform}              & Affine transform \\
\texttt{src.meta}                    & Metadata dict \\
\texttt{rasterio.plot.show(src)}    & Quick plot \\
\end{tabularx}

\sect{Coordinate Reference}
\begin{tabularx}{\linewidth}{@{}l@{\enspace}Y@{}}
\texttt{EPSG:4326}                   & WGS 84 (lon/lat) \\
\texttt{EPSG:3857}                   & Web Mercator \\
\texttt{pyproj.Transformer}         & CRS transforms \\
\texttt{t.transform(x, y)}          & Transform coords \\
\end{tabularx}

\end{multicols}

\end{document}
